% Generated draft: requires mathematical and source review.
\documentclass[11pt]{article}
\usepackage[margin=25mm]{geometry}
\usepackage{fontspec}
\setmainfont{Latin Modern Roman}
\usepackage{amsmath,amssymb,mathtools,unicode-math}
\setmathfont{Latin Modern Math}
\usepackage{enumitem,needspace,xurl,hyperref}
\setlength{\parindent}{0pt}
\setlength{\parskip}{5pt}
\setlength{\emergencystretch}{3em}
\newcommand{\N}{\mathbb{N}}
\newcommand{\Z}{\mathbb{Z}}
\newcommand{\Q}{\mathbb{Q}}
\newcommand{\R}{\mathbb{R}}
\newcommand{\C}{\mathbb{C}}
\newcommand{\F}{\mathbb{F}}
\newcommand{\A}{\mathbb{A}}
\newcommand{\PP}{\mathbb{P}}
\newcommand{\E}{\mathbb{E}}
\newcommand{\eps}{\varepsilon}
\newcommand{\dd}{\,\mathrm{d}}
\begin{document}
\section*{UnsolvedMath ID: 4400001}
\textbf{Source problem code:} catalogue entry
\subsection{Definitions and mathematical statement}
Give each finite alphabet the discrete topology and its bi-infinite sequence space the product topology. Define
\[
 X=\{0,1\}^{\Z},\qquad
 Y=\{y\in\{0,1,2\}^{\Z}:y_i\ne y_{i+1}\text{ for every }i\in\Z\}.
\]
On either space let $\sigma$ be the left shift, $(\sigma z)_i=z_{i+1}$. For a shift space $Z$ write
\[
 \operatorname{Per}(Z)=\bigcup_{k\ge1}\{z\in Z:\sigma^kz=z\},
 \qquad Z_{\mathrm{ap}}=Z\setminus\operatorname{Per}(Z),
\]
and give $Z_{\mathrm{ap}}$ the subspace topology. These complements are invariant under the shift.

Hochman's question is whether the two restricted dynamical systems are topologically conjugate: does there exist a homeomorphism
\[
 h:X_{\mathrm{ap}}\longrightarrow Y_{\mathrm{ap}}
 \qquad\text{such that}\qquad
 h\circ\sigma|_{X_{\mathrm{ap}}}=\sigma|_{Y_{\mathrm{ap}}}\circ h\ ?
\]
No extension of $h$ to the periodic points is required by this question.

The source states that $X$ and $Y$ are mixing shifts of finite type, both with entropy $\log2$, but are not isomorphic as the original shift systems. The research summary notes that their periodic points form countable dense sets and describes the complements as residual noncompact spaces. It repeatedly restates the target as mere homeomorphism of those spaces. That omits the shift-intertwining condition and is not the original question: the displayed target preserves topological conjugacy.
\subsection{Short English statement}
After deleting every periodic point from the full binary shift and the three-symbol shift forbidding equal neighbors, can one find a homeomorphism between the remaining spaces that commutes with the shift? The commuting requirement is essential; a plain homeomorphism of the underlying spaces does not answer the supplied statement.

The registry and research summary label the problem open, and the summary says no published resolution was found. Its search discussion focuses on homeomorphism rather than consistently on conjugacy, so it does not verify the original dynamical target's status. Provenance records no later authoritative resolution check beyond the 2010 source formulation.
\subsection{Sources}
\begin{itemize}
\item Local UnsolvedMath registry, permanent ID 4400001, source code catalogue entry, ``Pingree open problems --- Hochman problem 1,'' proposed by Mike Hochman, 2010. Dataset curation and metadata: Ulam AI and UnsolvedMath Contributors, CC BY 4.0.
\item \emph{Open problems from the 3rd Pingree Workshop on Dynamical Systems} (2010), Mike Hochman, Problem 1, \url{https://math.huji.ac.il/~mhochman/open-problems/pingree-open-problems.pdf}; pdf-text extraction, accessed 2026-07-29.
\item The supplied research summary provides a search-based OPEN-TRIAGE assessment but no later primary paper. Its plain-homeomorphism paraphrase differs from the source statement's topological-conjugacy target.
\item Source provenance marks status evidence NEEDS\_REVIEW and rights NEEDS\_REVIEW: the source presents the item as open, no later authoritative resolution check is recorded, and redistribution permission was not established despite verified public access. This qualification is separate from the dataset licence.
\end{itemize}
\end{document}
